\section{Pseudo-PDFs}

The involved  structure  of a quasi-PDF
$Q(y,P)$  can be attributed to the formal  fact that it is given by the Fourier $\nu$-transform  of
the function
${\cal M} (\nu, \nu^2/P^2)$,  in which    $\nu$ appears  both in the first and second
argument of the Ioffe-time distribution.
One should take   $P$-values that   are sufficiently large
to neglect the $\nu$-dependence coming from the second argument.

Another way  \cite{Radyushkin:2017cyf} is to try to {\it eliminate}  the \mbox{$z_3^2$-dependence}
induced by ${\cal M} (\nu, z_3^2)$. The main  idea is based
on the observation  that if one takes the
$\nu$-Fourier transform  of the modified function
${\cal M} (\nu, z_3^2)/D(z_3^2)$,  the $z_3\to 0$  limit
will   give the same PDF as the original Ioffe-time distribution,
provided that   $D(z_3^2)$ is a function of $z_3^2$ only
(but not of $\nu$) and is  equal to 1 for \mbox{$z_3^2=0$.}

Thus,  the strategy is to  find a function $D(z_3^2)$ whose \mbox{$z_3^2$-dependence }
would compensate, as much as possible,  the  $z_3^2$-dependence
of  ${\cal M} (\nu, z_3^2)$.
The next step is to fit the residual polynomial $z_3^2$-dependence
by polynomials of $z_3^2$ (they may be different for different values of $\nu$),
and in this way extrapolate the data to $z_3^2=0$ limit.
The Fourier transform
of  the resulting function  would correspond to
 the same PDF  as the $z_3^2$ limit of the original Ioffe-time  distribution
 ${\cal M} (\nu, z_3^2)$.


In the most lucky situation,  the ratio ${\cal M} (\nu, z_3^2)/D(z_3^2)$
would have no polynomial $z_3^2$-dependence  (or just a very mild one).
In particular, when
${\cal M} (\nu, z_3^2)$ factorizes, i.e.,
 ${\cal M} (\nu, z_3^2)= {\cal M} (\nu, 0) {\cal M} (0, z_3^2)$, one should take
  \mbox{$D(z_3^2) ={\cal M} (0, z_3^2)$}.  In this case,
 the reduced function
\begin{align}
{\mathfrak M} (\nu, z_3^2) \equiv \frac{ {\cal M} (\nu, z_3^2)}{{\cal M} (0, z_3^2)} \
 \label{redm}
\end{align}
is equal to $ {\cal M} (\nu, 0)$,
and
the  task  of obtaining  the $z_3\to 0$ limit is accomplished.

While there is no ``first principle'' reason for such a factorization,
one may expect that  the functions ${\cal M} (\nu, z_3^2)$
 for different $\nu$ have more or less similar dependence on $z_3$,
basically
reflecting the finite size of the nucleon.

As we mentioned already, the soft part of ${\cal M} (\nu, z_3^2)$  factorizes
if  the soft part of TMD ${\cal F} (x,k_\perp^2)$ factorizes.
That this happens, is
a standard  assumption of the TMD practitioners
(see, e.g., Ref.  \cite{Anselmino:2013lza}).
So, there are good chances that  this part of the
$z_3^2$-dependence of ${\cal M} (\nu, z_3^2)$
will be canceled or strongly reduced by the rest-frame function ${\cal M} (0, z_3^2) $.


On the lattice, there is another (and troublesome, see, e.g., Ref. \cite{Ishikawa:2016znu}) source
of $z_3$-dependence:
 the  $Z(z_3^2)$ factor
generated
by the   renormalization of
the gauge link  ${ \hat E} (0,z_3; A)$.
Fortunately, this  problematic     factor $Z(z_3^2)$
does not depend on $\nu$ and
is   the same for  the numerator and denominator
of the ratio  ${\mathfrak M} (\nu, z_3^2) $.
This provides another motivation for using ${\cal M} (0, z_3^2)$
as a factor $D(z_3^2)$.

 Thus, the proposal is to perform a lattice study
 of the reduced Ioffe-time  function ${\mathfrak M} (\nu, z_3^2)$.
 Even if it would have a residual polynomial $z_3^2$-dependence,
 it should  be much  easier to  extrapolate  this dependence to $z_3=0$,
 than the $z_3$-dependence of the original Ioffe-time  distribution
 $ {\cal M} (\nu, z_3^2)$.


Furthermore,   if one observes that  the ratio
$ {\mathfrak M} (\nu, z_3^2)$ does not have $z_3$-dependence,
one should  conclude that $ {\cal  M} (\nu, z_3^2)$ factorizes.
In fact, such a factorization has been already observed
several years ago
in the pioneering study \cite{Musch:2010ka} of the
transverse momentum distributions in  lattice QCD.


Still,  there is an unavoidable source of factorization breaking.
When $z_3$  is small,
$ {\cal M} (\nu, z_3^2)$ has   logarithmic $\ln z_3^2$ singularities
generating the perturbative  evolution of PDFs.
As we discussed,
 $1/z_3$
 is analogous then  to the renormalization parameter $\mu$
of  the scale-dependent PDFs $f(x,\mu^2)$ within the standard OPE approach.

    \begin{figure}[t]
   \centerline{ \includegraphics[width=4.0in]{images/MRER.pdf}  }
    \caption{Real  part   of model  distribution ${\cal M} (\nu)$  and
    the function $-B \otimes  {\rm Re} \,  {\cal M}$   that  governs its  evolution
    (the minus sign here is   for convenience of  placing  two curves on one   figure).
    \label{MBM}}
    \end{figure}

More specifically, for small values of $z_3$,  the  pseudo-PDF ${\cal P} (x,z_3^2)$
satisfies a  leading-order evolution equation
 with respect to $1/z_3$  that
is identical to  the evolution equation for $f (x,\mu^2)$ with respect to $\mu$.
The  evolution equation
for  the
reduced %
 Ioffe-time distribution ${\mathfrak M} (\nu, z_3^2)$ can also be written    \cite{Radyushkin:2017cyf}
    \begin{align}
    \frac{d}{d \ln z_3^2} \,
{\mathfrak M} (\nu, z_3^2)    &= - \frac{\alpha_s}{2\pi} \, C_F
\int_0^1  du \,   B ( u )   {\mathfrak M} (u \nu, z_3^2)
 ,
\label{EE}
 \end{align}
where  $C_F=4/3$, and the leading-order evolution kernel $B(u)$ for the non-singlet quark case
is given  \cite{Braun:1994jq} by
   \begin{align}
B (u)    &=  \left [ \frac{1+u^2}{1- u} \right ]_+  \
 ,
\label{Bu}
 \end{align}
 where $[ \ldots ]_+$  denotes the  ``plus'' prescription, i.e.
    \begin{align}
& \int_0^1  du \,    \left [ \frac{1+u^2}{1- u} \right ]_+     {\mathfrak M} (u \nu)
\nn & =
\int_0^1  du \,  \frac{1+u^2}{1- u} \,    [{\mathfrak M} ( \nu)-{\mathfrak  M} (u \nu)  ]\
 .
\label{plus}
 \end{align}


 Note that being a Fourier  transform,
  \begin{align}
  {\cal M} (\nu)  =
   \int_{-1}^1 dx \, f(x) \,
e^{ix\nu} \,  \  ,
 \label{Mf}
\end{align}
 the Ioffe-time distribution has real and imaginary parts even if the function $f(x)$
 is real (which  is the case with parton distributions). In particular,
   \begin{align}
  {\rm Re} \, {\cal M} (\nu)  =
   \int_{-1}^1 dx \, f(x) \,
\cos (x\nu)  \,  \  ,
 \label{ReMf}
\end{align}
 and
    \begin{align}
  {\rm Im} \, {\cal M} (\nu)  =
   \int_{-1}^1 dx \, f(x) \,
\sin (x\nu)  \,  \  .
 \label{ImMf}
\end{align}



  In \mbox{Fig. \ref{MBM}},  we show the function
  ${\rm Re}  \, {\cal M} (\nu)$ for a model PDF
      \begin{align}
q (x)= \frac{315}{32}  \sqrt{x} (1-x)^3 \theta ( 0\leq x \leq 1) \ .
\end{align}
Its integral is normalized to 1, and it  is nonzero for positive $x$ only,
  which corresponds  to  the absence of antiquarks.
As we shall see, this particular form    appears in the description of   actual  lattice data.
In \mbox{Fig. \ref{MBMI}}, we     show the  function
  ${\rm Im}  \, {\cal M} (\nu)$ for the same  model PDF.


    \begin{figure}[t]
     \centerline{ \includegraphics[width=3.8in]{images/MIEIkey.pdf}}
    \caption{Imaginary  part of model Ioffe-time distribution $  {\cal M} (\nu)$  and
    the function   $B \otimes  {\rm Im} \,  {\cal M}$  that governs its   evolution.
    \label{MBMI}}
    \end{figure}



 We also show in these figures
the convolution integrals governing the  evolution, namely
\mbox{$-B \otimes {\rm Re}\, {\cal M}(\nu)$}
and \mbox{$B \otimes {\rm Im}\, {\cal M}(\nu)$}.
The reader can notice that  $B \otimes {\rm }\,  {\cal M}(\nu)$  is zero for
$\nu=0$,  resulting from the vector current conservation.
As a consequence, the perturbative evolution leaves the rest-frame density ${\cal M} (0,z_3^2)$  (which is always real) unaffected.  In other words, the $\ln z_3^2$ terms
are present  only  in the numerator ${\cal M} ( \nu, z_3^2)  $
of the $ {\mathfrak M} (\nu, z_3^2)$ ratio,
but not in its ${\cal M} ( 0, z_3^2)  $  denominator.

Note also that  the evolution of the real part always leads
to a {\it decrease} of   ${{\rm Re} \, \cal M} ( \nu, z_3^2)$ when $z_3^2$ increases.
For the imaginary part, the evolution pattern is more complicated.
Namely,    below $\nu \sim 5.5$, the function  ${{\rm Im} \, \cal M} ( \nu, z_3^2)$
{\it increases}
 when $z_3^2$ increases. Only above $\nu \sim 5.5$, the evolution leads
 to a decrease of ${{\rm Im} \, \cal M} (u \nu, z_3^2)$ with $z_3^2$,
 and  the evolution pattern becomes  similar to that of the real part.
